\fsection{Ej 9 Pág 159}
\fsubsection
[\textcolor{waveBlue2}{Indicador clave: consumo de agua}]
{\textcolor{waveBlue2}{\boldit{Indicador clave: consumo de agua}}}
{
	En la elaboración del plan de sostenibilidad, la empresa puede tener como uno de sus objetivos
	principales la reducción del consumo de agua. Elabora un indicador clave que englobe ese objetivo
	y justifica tu respuesta, encuadrándola en tu sector productivo.
}

\begin{solution}

	\textbf{Indicador propuesto:}

	\textbf{Consumo de agua por unidad de producción} (m\textsuperscript{3}/unidad producida)

	\medskip

	\textbf{Definición:} Mide el volumen total de agua consumida en el proceso productivo dividido entre
	el número de unidades fabricadas o el valor de la producción obtenida en un período determinado.

	\medskip

	\textbf{Fórmula:}
	\[
		I_{\text{agua}} = \frac{\text{Volumen total de agua consumida (m}^3\text{)}}{\text{Unidades producidas}}
	\]

	\medskip

	\textbf{Justificación:}

	Este indicador permite medir de forma objetiva y comparable la eficiencia hídrica de la empresa a
	lo largo del tiempo. Al relativizar el consumo respecto a la producción, se eliminan las distorsiones
	causadas por variaciones en el volumen de actividad, lo que hace que el indicador sea válido tanto
	en períodos de alta como de baja producción.

	\medskip

	\textbf{Encuadre sectorial (sector industrial/manufacturero):}

	En el sector industrial, el agua se utiliza como materia prima auxiliar, en procesos de refrigeración
	y limpieza, y en la generación de vapor. Una reducción en este indicador refleja mejoras en la
	eficiencia de los procesos, como la implantación de circuitos cerrados de refrigeración, la
	reutilización de aguas grises o la sustitución de procesos húmedos por tecnologías en seco.

	\medskip

	\textbf{Meta orientativa:} Reducción del 10\,\% del indicador en un plazo de 2 años respecto al
	valor de referencia del ejercicio anterior.

\end{solution}
